% =====================================================================
% Abstract & Keywords
% =====================================================================
\begin{abstract}
In Taiwan, residents must hand their garbage to the collection truck in
person during a short stop, so missing the truck means waiting for the
next round. Existing services in Taichung City show truck positions and
estimated arrival times, but residents still have to decide on their own
which nearby collection point they can reach in time. We present
GarbageGo, a domain AI agent that turns this lookup task into a
time-feasibility recommendation. For each candidate point, the system
computes a slack, defined as the truck's estimated time of arrival minus
the walking time minus a safety buffer that grows with data staleness
and the number of remaining stops. It then ranks points by slack and
risk, offers alternatives, and reports when no feasible option exists.
A large language model interprets requests and explains results, while
every numerical value comes from dedicated GIS, ETA, and recommendation
engines. A 29-day measurement of Taichung's open vehicle API collected
980{,}775 deduplicated GPS points from 668 vehicles. The API publishes
snapshots about every five minutes, and 18\% of vehicles in one snapshot
had not reported for over 30 minutes. Of the license plates in the
collection-point dataset, 95.7\% appear in the vehicle data, covering
96.1\% of 20{,}090 points, and 95.6\% of point addresses were geocoded by
local address matching. These results show the data can support the
proposed design. We also define an evaluation plan that compares
nearest-point and slack-based strategies.
\end{abstract}

\begin{IEEEkeywords}
Internet of Things, open government data, estimated time of arrival,
LLM agent, waste collection, recommendation
\end{IEEEkeywords}
